% GENERATED FILE. Edit canonical lessons, then run npm run generate:books.
% canonical-source-hash: fnv1a64:a82048cb1cc819ff
% canonical-lessons: HI-C490-chala, HI-C490-taka, HI-C490-papri, HI-C490-kuredna, HI-C490-dag

\chapter{Cuts and Scrapes}
\label{ch:hi-cuts-and-scrapes}
\hlchaptermodality{490}

\begin{chapteropening}
\textbf{By the end of this chapter:} \emph{I can say how someone got a small injury.}

Complete the last lesson of chapter 490: I can say how someone got a small injury.
\end{chapteropening}

\section[\texorpdfstring{\hi{छाला} (chālā)}{छाला (chālā)}]{\hi{छाला} (chālā) — a blister from rubbing}

\label{lesson:HI-C490-chala}



\begin{quote}
Before the new one: say the Hindi for an ointment, then the Hindi for a herbal brew.
\end{quote}

\subsection*{You'll want to know: \hi{छाला}}
\textbf{\hi{छाला}} — \emph{chālā} — \textquotedblleft{}a blister from rubbing\textquotedblright{}.

\begin{quote}
\hi{नए} \hi{जूतों} \hi{से} \hi{मेरे} \hi{पैर} \hi{में} \hi{छाला} \hi{पड़} \hi{गया।}
\end{quote}

\begin{quote}
\emph{nae jūtoṁ se mere pair meṁ chālā paṛ gayā.}
\end{quote}

\begin{quote}
My new shoes gave me a blister on my foot.
\end{quote}

Masculine.

\begin{cousinweb}
\textbf{\hi{छाला}} (\emph{chālā}) is inherited: CDIAL lists Hindi \emph{chālā}, with the sense \textquotedblleft{}skin, hide\textquotedblright{}, under Sanskrit \emph{challi}, \textquotedblleft{}bark\textquotedblright{} (CDIAL 5005), beside \emph{chālī}, \textquotedblleft{}a small blister\textquotedblright{}. The \textquotedblleft{}blister\textquotedblright{} sense of \emph{chālā} itself, skin lifted from the flesh, is not given there.
\end{cousinweb}

\begin{grammarlens}[title={What you've built}]
One more word for describing what you lived and felt.
\end{grammarlens}

\subsection*{Your turn}
\begin{itemize}
\raggedright
  \item \emph{Say it:} \emph{chālā}
  \item \emph{Say it:} \emph{chālā}, once more
  \item \emph{Say it:} \emph{nae jūtoṁ se mere pair meṁ chālā paṛ gayā.} — My new shoes gave me a blister on my foot.
  \item \emph{Recall:} say the Hindi for an ointment, then the Hindi for a herbal brew, then say \emph{chālā} again
\end{itemize}

\begin{tcolorbox}[breakable,colback=teal!4,colframe=teal!35!black,title={Before you move on}]
What does \hi{छाला} mean? (\textquotedblleft{}A blister from rubbing\textquotedblright{}.) Say it once more.
\end{tcolorbox}
\section[\texorpdfstring{\hi{टाँका} (ṭā̃kā)}{टाँका (ṭā̃kā)}]{\hi{टाँका} (ṭā̃kā) — a stitch}

\label{lesson:HI-C490-taka}



\begin{quote}
Before the new one: say the Hindi for a herbal brew, then the Hindi for a blister from rubbing.
\end{quote}

\subsection*{You'll want to know: \hi{टाँका}}
\textbf{\hi{टाँका}} — \emph{ṭā̃kā} — \textquotedblleft{}a stitch\textquotedblright{}.

\begin{quote}
\hi{घाव} \hi{पर} \hi{चार} \hi{टाँके} \hi{लगे।}
\end{quote}

\begin{quote}
\emph{ghāv par cār ṭā̃ke lage.}
\end{quote}

\begin{quote}
The wound needed four stitches.
\end{quote}

Masculine.

\begin{cousinweb}
\textbf{\hi{टाँका}} (\emph{ṭā̃kā}) is inherited: CDIAL lists Hindi \emph{ṭā̃kā}, \textquotedblleft{}stitch, join\textquotedblright{}, under Sanskrit \emph{ṭáṅkati}, \textquotedblleft{}binds\textquotedblright{} (CDIAL 5432).
\end{cousinweb}

\begin{grammarlens}[title={What you've built}]
One more word for describing what you lived and felt.
\end{grammarlens}

\subsection*{Your turn}
\begin{itemize}
\raggedright
  \item \emph{Say it:} \emph{ṭā̃kā}
  \item \emph{Say it:} \emph{ṭā̃kā}, once more
  \item \emph{Say it:} \emph{ghāv par cār ṭā̃ke lage.} — The wound needed four stitches.
  \item \emph{Recall:} say the Hindi for a herbal brew, then the Hindi for a blister from rubbing, then say \emph{ṭā̃kā} again
\end{itemize}

\begin{tcolorbox}[breakable,colback=teal!4,colframe=teal!35!black,title={Before you move on}]
What does \hi{टाँका} mean? (\textquotedblleft{}A stitch\textquotedblright{}.) Say it once more.
\end{tcolorbox}
\section[\texorpdfstring{\hi{पपड़ी} (papṛī)}{पपड़ी (papṛī)}]{\hi{पपड़ी} (papṛī) — a scab; a thin crust}

\label{lesson:HI-C490-papri}



\begin{quote}
Before the new one: say the Hindi for a blister from rubbing, then the Hindi for a stitch.
\end{quote}

\subsection*{You'll want to know: \hi{पपड़ी}}
\textbf{\hi{पपड़ी}} — \emph{papṛī} — \textquotedblleft{}a scab; a thin crust\textquotedblright{}.

\begin{quote}
\hi{घुटने} \hi{के} \hi{घाव} \hi{पर} \hi{पपड़ी} \hi{जम} \hi{गई} \hi{है।}
\end{quote}

\begin{quote}
\emph{ghuṭne ke ghāv par papṛī jam gaī hai.}
\end{quote}

\begin{quote}
A scab has formed on the wound on my knee.
\end{quote}

Feminine.

\begin{cousinweb}
\textbf{\hi{पपड़ी}} (\emph{papṛī}) is a tadbhava: CDIAL lists Hindi \emph{papṛī}, \textquotedblleft{}thin crust, scab\textquotedblright{}, under Sanskrit \emph{párpaṭa} (CDIAL 7934). The thin, crisp \emph{papṛī} of the snack seller is the same word.
\end{cousinweb}

\begin{grammarlens}[title={What you've built}]
One more word for describing what you lived and felt.
\end{grammarlens}

\subsection*{Your turn}
\begin{itemize}
\raggedright
  \item \emph{Say it:} \emph{papṛī}
  \item \emph{Say it:} \emph{papṛī}, once more
  \item \emph{Say it:} \emph{ghuṭne ke ghāv par papṛī jam gaī hai.} — A scab has formed on the wound on my knee.
  \item \emph{Recall:} say the Hindi for a blister from rubbing, then the Hindi for a stitch, then say \emph{papṛī} again
\end{itemize}

\begin{tcolorbox}[breakable,colback=teal!4,colframe=teal!35!black,title={Before you move on}]
What does \hi{पपड़ी} mean? (\textquotedblleft{}A scab; a thin crust\textquotedblright{}.) Say it once more.
\end{tcolorbox}
\section[\texorpdfstring{\hi{कुरेदना} (kurednā)}{कुरेदना (kurednā)}]{\hi{कुरेदना} (kurednā) — to pick at; to scrape at}

\label{lesson:HI-C490-kuredna}



\begin{quote}
Before the new one: say the Hindi for a stitch, then the Hindi for a scab.
\end{quote}

\subsection*{You'll want to know: \hi{कुरेदना}}
\textbf{\hi{कुरेदना}} — \emph{kurednā} — \textquotedblleft{}to pick at; to scrape at\textquotedblright{}.

\begin{quote}
\hi{डाक्टर} \hi{ने} \hi{कहा} \hi{कि} \hi{पपड़ी} \hi{को} \hi{मत} \hi{कुरेदो।}
\end{quote}

\begin{quote}
\emph{ḍākṭar ne kahā ki papṛī ko mat kuredo.}
\end{quote}

\begin{quote}
The doctor said not to pick at the scab.
\end{quote}

Past, with \hi{ने} (he or she …): \hi{उसने} \hi{पपड़ी} \hi{कुरेदी}, \emph{usne papṛī kuredī}.

\begin{cousinweb}
\textbf{\hi{कुरेदना}} (\emph{kurednā}) is a tadbhava: CDIAL lists Hindi \emph{kurednā}, \textquotedblleft{}to scrape\textquotedblright{}, under \emph{*kura}, a form it reconstructs (CDIAL 3319). One can also \emph{kurednā} an old memory: stir it up again.
\end{cousinweb}

\begin{grammarlens}[title={What you've built}]
One more word for describing what you lived and felt.
\end{grammarlens}

\subsection*{Your turn}
\begin{itemize}
\raggedright
  \item \emph{Say it:} \emph{kurednā}
  \item \emph{Say it:} \emph{kurednā}, once more
  \item \emph{Say it:} \emph{ḍākṭar ne kahā ki papṛī ko mat kuredo.} — The doctor said not to pick at the scab.
  \item \emph{Recall:} say the Hindi for a stitch, then the Hindi for a scab, then say \emph{kurednā} again
\end{itemize}

\begin{tcolorbox}[breakable,colback=teal!4,colframe=teal!35!black,title={Before you move on}]
What does \hi{कुरेदना} mean? (\textquotedblleft{}To pick at; to scrape at\textquotedblright{}.) Say it once more.
\end{tcolorbox}
\section[\texorpdfstring{\hi{दाग़} (dāġ)}{दाग़ (dāġ)}]{\hi{दाग़} (dāġ) — a scar; a stain}

\label{lesson:HI-C490-dag}



\begin{quote}
Before the new one: say the Hindi for a scab, then the Hindi for to pick at.
\end{quote}

\subsection*{You'll want to know: \hi{दाग़}}
\textbf{\hi{दाग़}} — \emph{dāġ} — \textquotedblleft{}a scar; a stain\textquotedblright{}.

\begin{quote}
\hi{बचपन} \hi{की} \hi{चोट} \hi{का} \hi{दाग़} \hi{अरुण} \hi{के} \hi{माथे} \hi{पर} \hi{आज} \hi{भी} \hi{है।}
\end{quote}

\begin{quote}
\emph{bacpan kī coṭ kā dāġ Aruṇ ke māthe par āj bhī hai.}
\end{quote}

\begin{quote}
Arun still has the scar of a childhood injury on his forehead.
\end{quote}

Masculine.

\begin{cousinweb}
\textbf{\hi{दाग़}} (\emph{dāġ}) is usually traced to Persian \emph{dāġ}, \textquotedblleft{}a brand, a mark\textquotedblright{}: a stain on cloth or a scar on skin.
\end{cousinweb}

\begin{grammarlens}[title={What you've built}]
One more word for describing what you lived and felt.
\end{grammarlens}

\subsection*{Your turn}
\begin{itemize}
\raggedright
  \item \emph{Say it:} \emph{dāġ}
  \item \emph{Say it:} \emph{dāġ}, once more
  \item \emph{Say it:} \emph{bacpan kī coṭ kā dāġ Aruṇ ke māthe par āj bhī hai.} — Arun still has the scar of a childhood injury on his forehead.
  \item \emph{Recall:} say the Hindi for a scab, then the Hindi for to pick at, then say \emph{dāġ} again
\end{itemize}

\begin{tcolorbox}[breakable,colback=teal!4,colframe=teal!35!black,title={Before you move on}]
What does \hi{दाग़} mean? (\textquotedblleft{}A scar; a stain\textquotedblright{}.) Say it once more.
\end{tcolorbox}
